\documentclass[11pt]{article}
\usepackage[margin=1in]{geometry}
\usepackage{amsmath,amsthm,booktabs}
\usepackage[hidelinks]{hyperref}
\usepackage{xurl}
\setlength{\emergencystretch}{3em}
\newtheorem{theorem}{Theorem}
\newtheorem{corollary}[theorem]{Corollary}
\title{Mixed double branched covers of twenty-four A6 examples}
\author{Exact topology and existing-family scope}
\date{17 November 2026}
\begin{document}
\maketitle

The twenty-four knots with new A6 affine certificates have actual
double branched covers containing a hyperbolic piece. This excludes
both the graph-cover hypothesis of Boileau--Jang--Weidmann and
arbitrary rational replacements on a four-face BBK base.
These are precise family exclusions; worldwide named-equality
priority and comparison with other methods are separate questions.

\begin{theorem}\label{thm:outside}
For each knot in the following table, the actual double branched
cover $\Sigma_2(K)$ is closed, orientable and irreducible.
An incompressible separating torus cuts off a component whose
interior is homeomorphic to the listed finite-volume hyperbolic
census manifold:
\begin{center}\small
\begin{tabular}{lp{11.5cm}}
\toprule
component & knots\\
\midrule
m143 & $13a_j$ for $j=304,362,3232,3259$;\newline
$13n_j$ for $j=834,838,839,991,992,993,3814,3840$\\[2pt]
m239 & $13a_j$ for $j=623,629,3374,3375$;\newline
$13n_j$ for $j=1367,1370,1371,1374,1375,1376,3911,3914$\\
\bottomrule
\end{tabular}
\end{center}
In particular all twenty-four covers are nonhyperbolic and are not
graph manifolds.
\end{theorem}

\section{The literal diagrams and actual covering construction}

The exact input is \path{a6_cover_topology_inputs.json}.
It retains every literal PD, the independently packaged
Hoste--Thistlethwaite census PD and DT code, and a dart bijection
preserving the edge involution, cyclic crossing orders and
over/under distinction. All twenty-four selected maps preserve
orientation and over/under. The replay reads the current packaged
census under its pinned version and data-file hashes and verifies
these exact maps. Knot identity is not inferred from a homology
group, a numerical invariant or an unverified name attached to a PD.

For each PD, Spherogram builds the exterior with the diagram
meridian peripheral basis. The meridian coefficient-two orbifold
filling has a unique cyclic cover of degree two, and the primitive
lift of the branching slope has coordinates $(1,0)$.
Filling that lift gives the ordinary double cover branched over the
diagram knot. The resulting triangulation is converted to Regina
and simplified by bistellar moves.

The input pins the random seed, the exterior isomorphism signature
with peripheral data and the closed-cover signature.
Every replay reconstructs those exact signatures and checks validity,
connectedness, closedness and orientability.
The chosen closed triangulations have sixteen or seventeen tetrahedra.
Random choices select triangulations; a mismatched signature fails
rather than substituting a different unchecked triangulation.
No floating-point volume or numerical isometry test is used.

\section{Normal tori and certified hyperbolic pieces}

Regina's exact \texttt{isIrreducible()} predicate first checks the
closed triangulation. Vertex normal surfaces in quadrilateral
coordinates are then enumerated. Each input specifies a surface index
and the complete quadrilateral vector. The replay checks that exact
entry and requires it to be connected, orientable, two-sided and of
Euler characteristic zero. Thus it is an embedded torus.
The \path{NormalSurface.isIncompressible()} predicate returns true,
after its closed irreducible ambient-manifold hypotheses have been
verified.

Cutting along the surface gives exactly two components, each with one
torus boundary. The selected component is simplified, its boundary is
idealized, and it is simplified again. Its exact isomorphism signature
must occur in the m143 or m239 row of the author's certified cusped
census data. Both are five-tetrahedron census manifolds.
For example, the selected $13a_{304}$ component has signature
\texttt{fLLQcacdeeejkexjn}, an m143 entry.

The hyperbolicity import is Burton's Corollary~1.2 in the retained
author preprint \cite[printed p.~2]{burton}: the named census contains
no nonhyperbolic entries or omissions through eight tetrahedra.
The supporting \path{5-or.dat} lists exact minimal-triangulation
signatures; its retained copy is
\path{sources/burton-hypcensus-5-or.dat}.
Its bytes agree with the freshly downloaded author archive.
Equality of isomorphism signatures identifies exact face pairings.
This imports the certified census theorem at its stated scope;
the census-wide interval verification is not repeated here.

A finite-volume hyperbolic component separated by an incompressible
torus persists under refinement to the JSJ decomposition.
A closed irreducible graph manifold cannot contain such a component.
This is the standard consequence of the JSJ gluing criterion and
geometrization decomposition in
\cite[Proposition~1.6.2 and Theorem~1.7.6]{afw}.
The same incompressible torus excludes a closed hyperbolic structure.
These observations prove Theorem~\ref{thm:outside}.

\begin{corollary}\label{cor:bjw}
None of the twenty-four knots satisfies the graph-manifold
double-branched-cover hypothesis of the BJW18 theorem.
\end{corollary}
\begin{proof}
The published theorem explicitly assumes a graph-manifold double
branched cover \cite[Theorem~1.2]{bjw}.
Theorem~\ref{thm:outside} excludes that hypothesis.
\end{proof}

\begin{corollary}\label{cor:rational}
None of the twenty-four knots is obtained by arbitrary rational
replacements of the twist boxes of any sharp-four BBK base.
\end{corollary}
\begin{proof}
The retained canonical proof
\path{RATIONAL_REPLACEMENT_TOPOLOGY.tex}, Theorem~1
(label \texttt{thm:class}), states that an irreducible double
branched cover of a knot obtained by rational decoration of the medial
projection of a connected simple four-vertex plane graph is either
hyperbolic or a graph manifold.
A sharp-four BBK base has such a simple four-vertex dual.
Theorem~\ref{thm:outside} supplies an irreducible cover of neither type.
The rational replacements here may have arbitrary slopes, rotations,
mirrors and signs; preserving reflection labels is not required.
\end{proof}

The corollary concerns four-face bases.
The same canonical rational-replacement note proves that arbitrary
origins from higher-face bases are universal when rank and bridge
preservation are dropped. Such origins are not excluded.

\section{Exact replay and assurance}

The verifier is \path{verify_a6_cover_topology.py};
its receipt is \path{a6_cover_topology_verification.json}.
The environment is SnapPy 3.3.2, Spherogram 2.4.1,
\texttt{snappy\_15\_knots} 1.2.1,
\texttt{snappy\_manifolds} 1.4 and Regina 7.4.
The actual replay covers all twenty-four literal named diagrams,
closed triangulations, incompressible normal tori and exact cusped
components. Wrong knot names, changed hyperbolic census names and
corrupted normal vectors are rejected.

The combinatorial cover construction, filling and simplification
algorithms, Regina's recognition predicates and Burton's census
certificate are explicit trusted imports.
The incompressibility API requires a connected compact embedded
normal surface in a valid closed irreducible triangulation;
each condition is checked before invocation.
Cutting is performed by \texttt{cutAlong()}, not by a crushing
operation that might discard a component.

The separate A6 affine proof supplies the meridional and bridge
equalities; no branched-cover rank bound is a premise of this
topology theorem. Likewise the named-knot finite-group proof does not
require the topological family comparison.
These internal exact checks do not constitute formal verification,
independent human specialist review or exhaustive worldwide priority
verification. No periodic-parameter or other-knot generalization is
inferred from these twenty-four cases.

\begin{thebibliography}{9}
\bibitem{burton}
B. A. Burton, The cusped hyperbolic census is complete,
author preprint arXiv:1405.2695v1, Corollary~1.2, printed p.~2.
\url{https://arxiv.org/abs/1405.2695v1}.
Author data:
\url{https://people.smp.uq.edu.au/BenjaminBurton/code/hypcensus.zip}.
\bibitem{afw}
M. Aschenbrenner, S. Friedl and H. Wilton,
3-Manifold Groups, EMS (2015), final author manuscript,
Proposition~1.6.2 and Theorem~1.7.6, printed pp.~15--19.
\url{https://doi.org/10.4171/154}.
\bibitem{bjw}
M. Boileau, Y. Jang and R. Weidmann,
Meridional rank and bridge number for a class of links,
Pacific J. Math. 292 (2018), 61--84, Theorem~1.2.
\url{https://doi.org/10.2140/pjm.2018.292.61}.
\end{thebibliography}
\end{document}
